\documentclass{article}

% if you need to pass options to natbib, use, e.g.:
%     \PassOptionsToPackage{numbers, compress}{natbib}
% before loading neurips_2026

% The authors should use one of these tracks.
% Before accepting by the NeurIPS conference, select one of the options below.
% 0. "default" for submission
\usepackage{neurips_2026}
\usepackage{amsmath}
\usepackage{textcomp}
\usepackage{algorithm}     % Provides the floating "Algorithm" environment (optional)
\usepackage{algpseudocode} % Provides the "algorithmic" environment and command syntax



\usepackage{comment}
% the "default" option is equal to the "main" option, which is used for the Main Track with double-blind reviewing.
% 1. "main" option is used for the Main Track
%  \usepackage[main]{neurips_2026}
% 2. "position" option is used for the Position Paper Track
%  \usepackage[position]{neurips_2026}
% 3. "eandd" option is used for the Evaluations & Datasets Track
 % \usepackage[eandd]{neurips_2026}
 % if you want to opt-in for a de-anomymized submission (not recommended, but OK):
 %\usepackage[eandd, nonanonymous]{neurips_2026}
% 4. "creativeai" option is used for the Creative AI Track
%  \usepackage[creativeai]{neurips_2026}
% 5. "sglblindworkshop" option is used for the Workshop with single-blind reviewing
 % \usepackage[sglblindworkshop]{neurips_2026}
% 6. "dblblindworkshop" option is used for the Workshop with double-blind reviewing
%  \usepackage[dblblindworkshop]{neurips_2026}

% After being accepted, the authors should add "final" behind the track to compile a camera-ready version.
% 1. Main Track
 % \usepackage[main, final]{neurips_2026}
% 2. Position Paper Track
%  \usepackage[position, final]{neurips_2026}
% 3. Evaluations & Datasets Track
 % \usepackage[eandd, final]{neurips_2026}
% 4. Creative AI Track
%  \usepackage[creativeai, final]{neurips_2026}
% 5. Workshop with single-blind reviewing
%  \usepackage[sglblindworkshop, final]{neurips_2026}
% 6. Workshop with double-blind reviewing
%  \usepackage[dblblindworkshop, final]{neurips_2026}
% Note. For the workshop paper template, both \title{} and \workshoptitle{} are required, with the former indicating the paper title shown in the title and the latter indicating the workshop title displayed in the footnote.
% For workshops (5., 6.), the authors should add the name of the workshop, "\workshoptitle" command is used to set the workshop title.
% \workshoptitle{WORKSHOP TITLE}

% "preprint" option is used for arXiv or other preprint submissions
 % \usepackage[preprint]{neurips_2026}

% to avoid loading the natbib package, add option nonatbib:
%    \usepackage[nonatbib]{neurips_2026}

\usepackage[utf8]{inputenc} % allow utf-8 input
\usepackage[T1]{fontenc}    % use 8-bit T1 fonts
\usepackage{hyperref}       % hyperlinks
\usepackage{url}            % simple URL typesetting
\usepackage{booktabs}       % professional-quality tables
\usepackage{amsfonts}       % blackboard math symbols
\usepackage{nicefrac}       % compact symbols for 1/2, etc.
\usepackage{microtype}      % microtypography
\usepackage{xcolor}         % colors
\usepackage{graphicx}       % figures

% Note. For the workshop paper template, both \title{} and \workshoptitle{} are required, with the former indicating the paper title shown in the title and the latter indicating the workshop title displayed in the footnote. 
\title{Convincing Biology with Language: Offline Learning of Prompt-Conditioned Biological Interventions}


% The \author macro works with any number of authors. There are two commands
% used to separate the names and addresses of multiple authors: \And and \AND.
%
% Using \And between authors leaves it to LaTeX to determine where to break the
% lines. Using \AND forces a line break at that point. So, if LaTeX puts 3 of 4
% authors names on the first line, and the last on the second line, try using
% \AND instead of \And before the third author name.


\author{%
  Nam Le \\
  Department of Computer Science \\
  University of Vermont \\
  Burlington, VT 05405 \\
  \texttt{nam.le@uvm.edu}
  \And
  Michael Levin \\
  Allen Discovery Center \\
  Tufts University \\
  Medford, MA 02155 \\
  \texttt{michael.levin@tufts.edu}
  \And
  Josh Bongard \\
  Department of Computer Science \\
  University of Vermont \\
  Burlington, VT 05405 \\
  \texttt{jbongard@uvm.edu}
}


\begin{document}


\maketitle


\begin{abstract}
Artificial intelligence has transformed how people interact with complex technical systems across domains -- from writing code to designing molecules -- enabling non-experts to accomplish sophisticated tasks through natural language prompts. We envision an analogous future for biotechnology: scientists and practitioners directing living matter by describing what they want, with learned models translating those descriptions into the interventions that convince cells, organoids, or engineered biological systems to behave accordingly. Learning this mapping requires paired language-intervention-outcome data, which is expensive to collect at scale. Recent work has demonstrated language-conditioned policies for simulated cellular control, but these rely on hand-engineered reward functions and task-specific architectures, and do not exploit existing observational archives, requiring fresh experiments for every new objective. Here we demonstrate, for the first time, that a policy network can be trained offline, from a fixed archive of intervention-outcome pairs, to translate natural language prompts into interventions shown to reliably elicit desired behaviors in \textit{in vitro} motile organoids (Xenobots). We frame this as offline contextual-bandit learning, with a three-stage pipeline: (1) a prompt-to-intervention (P2I) network maps language embeddings to intervention parameters, (2) a video archive lookup retrieves the historical Xenobot response to that intervention, and (3) a pretrained vision-language model provides a zero-shot reward signal scoring how well the observed behavior matches the original prompt. Optimizing the P2I network end-to-end with CMA-ES, and resolving objective interference between opposing behavioral objectives via a multi-head architecture, we demonstrate generalization to 40 unseen test prompts and 44 previously unobserved Xenobot videos, with 6 of 8 phenotypes exceeding the random-retrieval baseline. This work shows that language-guided biological control is learnable from purely observational data, without online experiments or hand-crafted rewards, and constitutes an intermediate step toward convincing biology with natural language.


\end{abstract}
% ---------------------------------------------------------------------------------------
\section{Introduction}

%Para 1: The vision (natural language → biological intervention), kept to ~4 sentences

Controlling living systems is a long-standing goal across biology, medicine, and synthetic biology, yet the interface remains fundamentally low-level: practitioners specify interventions parameter-by-parameter, through trial-and-error informed by domain expertise. In other technical domains, natural language has emerged as a high-level interface to complex systems: users now specify goals — generate this code, design this molecule, plan this trajectory — and learned models translate intent into low-level action. This raises a question central to applied ML: can a similar interface be learned for biological systems, mapping natural-language descriptions of desired behavior to physical interventions that elicit it?


%Para 2: Why it's hard (paired data expensive; conflicting objectives cause gradient interference; no formal grounding from language to bioelectric stimulus)

Two recent threads make this question tractable. Language-conditioned policies have been shown to control simulated cellular systems~\cite{Le2025zapgpt, le2025giving}, and targeted physical interventions — electrical, optogenetic, chemical — have been shown to elicit reproducible behaviors in real cells and tissues~\cite{zajdel2020scheepdog, levin2021bioelectric}. However, extending this to real organisms under open-ended language faces two coupled barriers. First, collecting paired (prompt, intervention, outcome) data online is prohibitively expensive: each datum is a wet-lab experiment. Second, the objective space is inherently multi-task and contains opposing goals — \textit{e.g.}, ``stop'' versus ``accelerate'' — which create objective interference when optimized through a shared output~\cite{pcgrad, chen2018gradnorm}. Existing approaches sidestep both barriers by restricting to hand-crafted reward functions or single-task architectures, and as a result do not generalize to open-ended natural language.


%Para 3: Our insight (archives already contain intervention records; VLM can score alignment; multi-head architecture decouples conflicting objectives)

We address both barriers by reframing language-to-intervention as an offline learning problem. Existing biological archives already contain large numbers of intervention-outcome pairs collected for purposes other than language-conditioned learning; we treat these as a fixed dataset over which a policy can be trained without any new experiments. To supply a reward signal in the absence of labeled prompt-outcome alignment, we use a pretrained vision-language model (VLM) as a zero-shot scorer~\cite{rocamonde2023vision}, evaluating how well an archived outcome matches a natural-language description. To address the second barrier, the policy uses a shared backbone with one output head per behavioral phenotype, enabling independent optimization across opposing objectives while preserving shared semantic representation~\cite{ruder2017overview}.


%Para 4: Contributions (3 bullets, concrete and measurable)
Our contributions are:
(1) A demonstration that the language-to-intervention mapping can be learned offline, from a fixed archive of intervention-outcome pairs collected without language annotation, eliminating the need for new wet-lab experiments per objective;
(2) An end-to-end pipeline combining a prompt-to-intervention policy network, archive lookup, and VLM-based zero-shot rewards, optimized with CMA-ES; and
(3) A multi-head policy architecture that resolves objective interference between opposing behavioral objectives, enabling a single network to fit antagonistic prompts (e.g.\ ``stop'' and ``accelerate'') simultaneously.

% ---------------------------------------------------------------------------------------
\section{Related Work}
\label{sec:related}

Large language models have demonstrated remarkable ability to serve as general-purpose interfaces to complex technical domains. Scaling has enabled few-shot task learning~\cite{brown2020language}, with subsequent work showing LLMs can generate functional code~\cite{chen2021evaluating}, design molecules~\cite{jumper2021highly}, and assist in scientific discovery~\cite{zheng2025large}. The common pattern is that users specify goals through natural language rather than low-level parameters. However, applying LLMs to biological control has proven difficult. Most prior work in this direction restricts to simulation, where prompts map to cell parameters with deterministic outcomes~\cite{Le2025zapgpt, le2025giving}. Direct application to living organisms is rare, as the mapping from language to biological intervention lacks formal grounding -- unlike code, biology offers no closed-form semantics for what move faster should mean as an electrical stimulus. Moreover, collecting paired data (prompt, intervention, outcome) requires expensive real-time experiments.



%Para 2: Synthetic biology / Xenobots / bioelectric control

Synthetic biology has produced sophisticated engineered systems with autonomous behavior~\cite{kriegman2021kinematic, blackiston2021cellular}. Xenobots, living robots assembled from \textit{Xenopus laevis} cells, exhibit spontaneous coordinated motion and adaptive morphology~\cite{kriegman2021kinematic}. Bioelectric networks encode high-level anatomical information, suggesting goal-directed control mechanisms~\cite{levin2021bioelectric}. Yet control of these systems remains manual and low-level: practitioners tune stimulation parameters (duration, current, frequency) through trial-and-error, without any high-level behavioral interface, and the intervention-to-behavior mapping is non-linear and poorly characterized. Targeted physical interventions are nonetheless effective at the parameter level: dynamic electric cues can steer cell migration with precision~\cite{zajdel2020scheepdog}, and optogenetic stimulation can probe biological decision-making~\cite{ebrahimkhani2021synthetic}. What is missing is a learned interface between high-level behavioral specifications and these low-level intervention protocols.



%Para 3: VLMs as reward models + offline RL + multi-task gradient conflicts

Bridging language and biological behavior requires two innovations. Vision-language models trained on joint image-text embeddings~\cite{radford2021learning, li2022blip} can be used directly as zero-shot reward models, scoring the alignment between an observation and a language description~\cite{rocamonde2023vision, Le2025zapgpt}. This approach has been explored in robotics~\cite{liang2023code, wang2024rl}; biological applications remain limited because the relevant behaviors are subtle, the visual backgrounds cluttered, and the outcomes stochastic. Learning a single policy across opposing behavioral goals (``stop'' vs.\ ``accelerate'') is a multi-task problem, and shared outputs are known to suffer objective interference~\cite{pcgrad, chen2018gradnorm}; task-specific heads over a shared backbone are a standard architectural remedy~\cite{wang2022can, ruder2017overview}. 

Finally, training from a fixed archive of intervention-outcome pairs places this work in the offline RL setting~\cite{levine2020offline}, and specifically in the contextual-bandit regime: there are no state transitions, only terminal rewards. This regime is well-suited to derivative-free black-box optimization, for which CMA-ES~\cite{hansen2001completely} is a standard choice — gradient-free, robust to multi-modal landscapes, and effective in the parameter dimensionalities considered here.



% ---------------------------------------------------------------------------------------
\section{Methods}
\label{sec:methods}

%Our approach (Fig.~\ref{fig:overview}) is as follows. A natural language prompt $p$ (Sect.~\ref{prompts}) is passed to an AI model that transforms it into an intervention $i$ (Sect.~\ref{p2i}). The archive (Sect.~\ref{archive}), partitioned into training and test splits (Sect.~\ref{sec:split}), is searched for the intervention $i'$ closest to $i$. The video $v$ of how an organoid reacted to $i'$ is extracted (Sect.~\ref{videos}) and sent, along with $p$, to a VLM that returns a score $s$ denoting how well the organoid in $v$ ``obeys'' $p$ (Sect.~\ref{vlm}). The backward pass tunes the policy to produce higher $s$ over a set of training prompts $P$ (Sect.~\ref{sec:training}).

Our approach (Fig.~\ref{fig:overview}) has two phases. In the forward pass, a natural-language prompt (Sect.~\ref{prompts}) is mapped by a policy to an intervention (Sect.~\ref{p2i}); the closest intervention is retrieved from an archive (Sect.~\ref{archive}), which has been split into training and test interventions (Sect.~\ref{sec:split}); the resulting video (Sect.~\ref{videos}) and prompt are scored by a VLM (Sect.~\ref{vlm}). In the backward pass, evolutionary search optimizes the policy to maximize VLM scores on training prompts (Sect.~\ref{sec:training}).

\begin{figure}[!t]
\centering
\includegraphics[width=1\textwidth]{figures/pipeline.pdf}
\caption{The pipeline. (A) User prompt $\to$ (B) P2I network $\to$ intervention duration $\to$ (C) archive lookup $\to$ (D) pre/post motion heatmaps $\to$ (E) VLM score. (F) CMA-ES evolves the P2I population on training prompts; (G) the trained P2I generalizes to held-out prompts and batches. Detailed in Sect.~\ref{sec:methods}.}
\label{fig:overview}
\end{figure}

\subsection{Prompts.}
\label{prompts}
Users specify desired behavior through natural language descriptions (e.g., ``make Xenobots stop moving'', ``increase activity''). Each prompt is embedded using a pretrained sentence transformer (SentenceBERT; \cite{reimers2019sentence}), producing a 384-dimensional embedding that captures semantic meaning. This embedding serves as the input to the P2I network.

\subsection{Prompt-to-Intervention (P2I):}
\label{p2i}

The policy network $\pi_\Theta$ maps prompt embeddings to behavior-conditioned intervention durations. Given prompt embedding $\mathbf{e} = \phi(p) \in \mathbb{R}^{384}$ (SentenceBERT), a shared convolutional backbone extracts features (detailed architecture and parameter counts in Appendix A):
\begin{equation}
\mathbf{f} = \text{GlobalAvgPool}(\text{Conv1D}_{384 \to 16}(\mathbf{e}))
\end{equation}
where the three-layer CNN reduces dimensionality ($384 \to 64 \to 32 \to 16$) with ReLU activations and dropout (0.2).

We compare two output designs. The \textit{single-head baseline} produces one scalar for all behaviors:
\begin{equation}
\mu_{\text{single}} = \sigma(\text{MLP}_{16 \to 1}(\mathbf{f}; \Theta_{\text{single}}))
\end{equation}

Under this design, prompts that demand opposing interventions -- ``stop'' (short stimulus) versus ``accelerate'' (long stimulus) -- backpropagate conflicting signals through the same output, so the network is pulled toward incompatible targets. Our \textit{multi-head design} decouples this by giving each of the 8 behaviors its own output head:
\begin{equation}
\mu_b(\mathbf{e}) = \sigma(\text{MLP}_{16 \to 32 \to 1}(\mathbf{f}; \Theta_b)), \quad b \in \{0, 1, \ldots, 7\}
\end{equation}

Independent heads allow specialization (each behavior optimizes its $\Theta_b$ independently) while the shared backbone (hard parameter sharing~\cite{ruder2017overview}) encodes semantic features useful across all behaviors. With a 7,984-parameter backbone and 577 parameters per behavior head, the 8-behavior multi-head architecture uses 12,600 total parameters versus 8,561 for single-head, a 47.2\% increase that enables behavior specialization. During training, prompt-behavior pairs are routed to the correct head via oracle labels; at test time, routing uses cosine similarity to the nearest training-prompt embedding, enabling generalization to unseen phrasings.

\subsection{The Archive.}
\label{archive}

Our dataset comprises 143 Xenobot intervention videos collected from an academic institution's bioelectricity research center
(between October 2025 and January 2026. Videos available upon request with institutional approval). Each video captures pre- and post-intervention 
behavior at 30 fps, 1280×960 resolution. Xenobot responses depend on multiple intervention parameters (current magnitude, stimulation angle, frequency, duration); we fix current, angle, and frequency, and treat \textit{stimulus duration} $d \in [1, 327]$ seconds (median: 131s) as the single intervention variable. Each batch consists of (1) raw video, (2) a command log with exact intervention parameters and timing, and (3) trajectory files. Restricting variation to a single parameter isolates the language-to-intervention mapping from confounds in the underlying biology.

\subsection{Data Splitting}
\label{sec:split}

The 143-batch archive is stratified into 99 training and 44 test batches, balancing difficulty tiers within each split. This stratification exploits a natural property of the data: VLM baseline scores (computed on random batch selection) reveal a bimodal distribution. \textit{Easy-tier} behaviors (stop, motion\_reduction, elimination, confinement, response\_magnitude) receive consistently high baseline scores averaging 0.74, indicating the archive contains abundant examples of these phenotypes. \textit{Hard-tier} behaviors (directedness, vigor, motion\_increase) receive sparse, low baseline scores averaging 0.16, indicating these phenotypes are rare in the archive (Table~\ref{tab:behaviors}). This distribution justifies the stratified train-test split and defines the difficulty landscape: behaviors that are naturally abundant require minimal prompting, while rare behaviors demand genuine language-conditioned learning to discover them.

We use 80 total prompts: 40 for training and 40 held-out for testing. The 40 training prompts (8 behaviors $\times$ 5 prompts) are partitioned 3:2 per behavior into optimization and validation prompts. The 99 training batches support both optimization and validation; the 44 test batches remain held-out.

\subsection{Videos.}
\label{videos}

Raw Xenobot videos (30 fps, 1280×960) contain optical noise from bubbles and electrode reflections. To extract clean behavioral signals, we denoise in three steps: (1) frame differencing to localize motion regions, (2) HSV saturation thresholding ($s > 20$) to reject low-saturation artifacts (bubbles, reflections), preserving pigmented Xenobot regions, and (3) morphological cleanup of the resulting mask (detailed in Appendix D).

Dense optical flow (Farneback method) computed over the denoised mask yields per-frame motion magnitude. We accumulate motion magnitude separately over pre-intervention and post-intervention periods, producing two motion heatmaps that visualize behavior change:

\begin{equation}
\label{eq:delta_h}
\Delta H = \sum_t \text{Magnitude}_{\text{masked}}^{\text{post}}(t) - \sum_t \text{Magnitude}_{\text{masked}}^{\text{pre}}(t)
\end{equation}

Figure~\ref{fig:exemplars_fulldata} shows examples of pre- and post-intervention motion heatmaps computed from four distinct batches in the archive. These heatmaps serve as visual inputs to the VLM reward signal (Sect.~\ref{vlm}), which evaluates alignment between the observed behavioral change and the command prompt given to the P2I network.


\begin{figure}[!t]
% Nam's version: https://docs.google.com/presentation/d/1Bzna4xfTwl7mqqzZkNj0ef2vDcF9dkwyrHuAkI6o8DM/edit?slide=id.g3dde17d994f_0_0#slide=id.g3dde17d994f_0_0
% Josh's version: https://docs.google.com/presentation/d/1paMzR_P6F8OOklS7JMjOY_RqZWeoEXUlWE9KKPiYvlc/edit?usp=sharing
\centering
% Nam's version
%\includegraphics[width=\textwidth]{figures/full_figures_2x2.png}
% Josh's version
\centerline{\includegraphics[width=0.5\textwidth]{figures/Fig2a.pdf}
            \includegraphics[width=0.5\textwidth]{figures/Fig2b.pdf}}
\centerline{\includegraphics[width=0.5\textwidth]{figures/Fig2c.pdf}
            \includegraphics[width=0.5\textwidth]{figures/Fig2d.pdf}}
%\caption{\textbf{Four examples of language-guided behavior produced by the trained P2I network.} For each prompt, the network outputs an intervention duration $i$, rounded to the nearest archived duration $i'$; panels show the pre- (blue) and post-intervention  (red) Xenobot trajectories in the matched archive batch. Yellow dashed circles denote the chamber boundary (4 mm diameter; 142 pixels/mm). Panel (d) is notable: prompted with ``Make the bot disappear,'' the network selects a long stimulation duration ($i' = 225$\,s) that retrieves a batch in which the Xenobot was most likely destroyed by prolonged electrical stimulation. The system thus recovers, from purely observational data, the relationship between long-duration stimulation and Xenobot destruction --- and correctly maps a natural-language description of that outcome onto the corresponding intervention.}
\caption{\textbf{Four examples of language-guided behavior produced by the trained P2I network.}
For each prompt, the network outputs an intervention duration $i$, rounded to the 
nearest archived duration $i'$; panels show the pre- (blue) and post-intervention 
(red) Xenobot trajectories in the matched archive batch. Yellow dashed circles 
denote the chamber boundary (4 mm diameter; 142 pixels/mm).}
\label{fig:exemplars_fulldata}
\end{figure}

\subsection{The VLM.}
\label{vlm}

For each (video, prompt) pair, we compute an alignment score using a frozen vision-language model (Qwen3.5-397B served via Ollama Cloud). The VLM takes as input: (1) pre- and post-intervention motion heatmaps from Eq.~\ref{eq:delta_h}, and (2) a query template populated with the command prompt given to the P2I network: ``Compare pre (blue) vs post (red) heatmaps. How well does post-intervention behavior match: \{prompt\}? Score 0 - 1.'' The VLM's scalar output $s \in [0, 1]$ represents zero-shot alignment - how faithfully the observed behavior change matches the requested prompt -- and serves as the reward signal. Because this score is deterministic per (video, prompt, behavior) triple under greedy decoding, we precompute the full score tensor $\mathbf{S} \in \mathbb{R}^{143 \times 80 \times 8}$ once, prior to optimization, eliminating redundant VLM calls during training. Detailed VLM query templates are provided in Appendix A.

\subsection{Training}
\label{sec:training}

We optimize policy parameters $\theta$ to maximize average reward across all behaviors and their training prompts:
\begin{equation}
\max_\theta \frac{1}{8} \sum_{b=1}^{8} \mathbb{E}_{\mathbf{e} \in P_b}[r(\theta, \mathbf{e}, b)]
\label{eq:objective}
\end{equation}
where $b \in \{1, \ldots, 8\}$ is the behavior index, $P_b$ is the set of training prompts for behavior $b$, the action $a \in [0, 1]$ is the intervention duration produced by the policy on prompt embedding $\mathbf{e} \in \mathbb{R}^{384}$, and the reward $r \in [0, 1]$ is the VLM score (Sect.~\ref{vlm}). Each behavior's reward is averaged across its training prompts. This is a one-step contextual bandit: terminal rewards only, no state transitions, 143 samples per behavior~\cite{may2012optimistic, wang2017optimal}. We adopt CMA-ES~\cite{hansen2001completely}, a derivative-free algorithm suited to sparse, non-differentiable rewards, with population size $\lambda = 4 + \lfloor 3 \ln(n) \rfloor = 32$ for $n = 12{,}600$ multi-head parameters, run for 100 generations with initial mutation strength $\sigma_0 = 0.3$.

% Josh's replacement Table 1 ----------------------------------
\begin{table}[h]
\centering
\footnotesize
\caption{The eight behavioral phenotypes used in this work. Operational definitions characterize the post-intervention motion pattern that defines each behavior; example prompts illustrate the natural-language descriptions used at training and test time. Train/Test baselines report mean VLM scores under random archive matching; these represent the inherent average behavior score of the dataset, achieved by randomly sampling durations.}
\label{tab:behaviors}
\begin{tabular}{p{2.2cm}p{2.4cm}cc}
\toprule
Behavior & Example prompt & Train & Test \\
\midrule
\multicolumn{4}{l}{\textit{Decrease in motion}} \\
\midrule
\texttt{stop} & ``Stop moving.'' & 0.750 & 0.743 \\
\texttt{motion\_reduction} & ``Move slower.'' & 0.763 & 0.810 \\
\midrule
\multicolumn{4}{l}{\textit{Increase in motion}} \\
\midrule
\texttt{motion\_increase} & ``Move faster.'' & 0.140 & 0.135 \\
\texttt{vigor} & ``Show vigor.'' & 0.156 & 0.152 \\
\texttt{response\_ magnitude} & ``Dramatic behavior change.'' & 0.762 & 0.824 \\
\midrule
\multicolumn{4}{l}{\textit{Structure of motion}} \\
\midrule
\texttt{directedness} & ``Show purposeful movement.'' & 0.150 & 0.190 \\
\midrule
\multicolumn{4}{l}{\textit{Spatial extent}} \\
\midrule
\texttt{confinement} & ``Stay confined to center.'' & 0.631 & 0.654 \\
\texttt{elimination} & ``Bot ceases all activity.'' & 0.700 & 0.758 \\
\bottomrule
\end{tabular}
\end{table}
% ---------------------------------------------------------------------------------------

\section{Results}
\label{sec:results}

%We first report training dynamics for the multi-head CMA-ES policy and the single-head baseline (Sect.~\ref{sec:results-training}), then evaluate the trained multi-head policy's generalization to held-out prompts and videos (Sect.~\ref{sec:testing}).


\subsection{Training}
\label{sec:results-training}

We compared CMA-ES variants (Table~\ref{tab:training_overall}): single-head and multi-head architectures, against random archive sampling as a reference control. The random baseline samples uniformly random durations from the archive for each prompt; over many samples, this converges to the dataset's inherent average behavior score (0.510 across all behaviors).

CMA-ES multi-head achieved $0.9269 \pm 0.0096$ training fitness and $0.7522 \pm 0.0137$ validation fitness, demonstrating strong learning and generalization. Single-head CMA-ES reached $0.6664 \pm 0.0218$ training fitness but only $0.3999 \pm 0.1124$ validation fitness, indicating severe overfitting: the shared backbone struggles to balance opposing objectives, causing poor transfer to held-out data. Multi-head specialization enables strong training and generalization despite the additional parameters.

\begin{table}[h]
\centering
\caption{Overall training and validation fitness across optimization methods (mean $\pm$ SD across 30 seeds). Baseline is the mean VLM score from the training dataset. Values in parentheses show improvement (+) or degradation (-) vs. baseline.}
\label{tab:training_overall}
\begin{tabular}{lcc}
\toprule
Method & Train Fitness & Val Fitness \\
\midrule
CMA-ES (multi-head) & $0.9269 \pm 0.0096$ (+81.5\%) & $0.7522 \pm 0.0137$ (+47.5\%) \\
CMA-ES (single-head) & $0.6664 \pm 0.0218$ (+30.6\%) & $0.3999 \pm 0.1124$ (-21.6\%) \\
Baseline & $0.510$ & $0.510$ \\
\bottomrule
\end{tabular}
\end{table}




% Nam is here: I am rewriting this -- I will report a table for overall score for GT1, GT2 vs baseline (aggreate for 8 behaviors), showing significance. Tables for each behavior will be on appendix. We may keep the figure for each behavior, which makes more visual, and no need to repeat same things.

\subsection{Testing}
\label{sec:testing}

The training validation results revealed that both single-head CMA-ES and REINFORCE fail to learn across the full behavioral repertoire: single-head suffers from severe overfitting ($0.3999$ validation fitness, -21.6\% vs baseline), while REINFORCE achieves only marginal gains that do not transfer ($0.426$ validation fitness, -16.5\% vs baseline). Although these methods may reach acceptable scores on individual easy behaviors, they cannot achieve a compromise solution balancing all 8 objectives. Accordingly, we focus generalization analysis on the successful approach: CMA-ES multi-head, which demonstrated strong validation performance ($0.7522$, +47.5\% vs baseline). Detailed per-behavior analysis of single-head and REINFORCE is provided in Appendix C.

We evaluate generalization across two protocols. \textbf{GT1} (Generalization to novel prompts) pairs 40 held-out test prompts with the 99 training-set videos, measuring transfer to unseen language. \textbf{GT2} (Generalization to novel batches) pairs the same 40 test prompts with 44 held-out videos from a different experimental cohort, testing transfer across both prompt and intervention space.

First, we compute overall average test scores across all 8 behaviors and 30 seeds (Table~\ref{tab:test_overall}). Both conditions show highly significant improvement over the random baseline. GT1 achieves strong transfer to novel prompts ($0.866$ mean score, $+70.9\%$ improvement, $p < 10^{-41}$). GT2 shows substantial degradation when batches are also novel ($0.627$ mean score, $+23.8\%$ improvement, $p < 10^{-20}$), but remains significantly above baseline.

\begin{table}[h]
\centering
\caption{Overall test generalization performance: averaged across all 8 behaviors and 30 seeds. Baseline is the mean VLM score from the dataset (average of all possible random intervention selections). GT1 demonstrates strong transfer to novel prompts; GT2 shows moderate transfer when both prompts and batches are novel.}
\label{tab:test_overall}
\small
\begin{tabular}{lccc}
\toprule
Condition & Mean Score & Improvement & $p$-value \\
\midrule
Baseline & 0.510 & — & — \\
GT1 (test prompts × train videos) & 0.866 & $+70.9\%$ & $p < 10^{-41}$ \\
GT2 (test prompts × test videos) & 0.627 & $+23.8\%$ & $p < 10^{-20}$ \\
\bottomrule
\end{tabular}
\end{table}

Per-behavior results reveal distinct generalization patterns (Appendix B). Under GT1, all eight behaviors exceed $p < 10^{-7}$, with gains scaling inversely to baseline difficulty: the three hardest behaviors -- \texttt{motion\_increase} (baseline 0.140), \texttt{directedness} (0.150), and \texttt{vigor} (0.156) --- achieve $+425\%$, $+388\%$, and $+375\%$ improvements respectively, reaching 0.73--0.74. Easy-tier behaviors near ceiling show smaller relative but substantial absolute gains, with \texttt{response\_magnitude} reaching 0.980.

Under GT2, six behaviors maintain positive transfer. The opposing behaviors \texttt{vigor} (+307\%), \texttt{motion\_increase} (+186\%), and \texttt{directedness} (+128\%) remain substantially improved. However, \texttt{confinement} and \texttt{motion\_reduction} show negative transfer (-53.1\% and -26.9\% respectively), likely reflecting test-cohort duration distribution differences that affect these phenotypes' sensitivity to stimulus duration. We return to this point in Sect.~\ref{sec:Discussion}. Further details, including per-behavior statistics (Cohen's d, R²) and inter-behavior correlations, are provided in Appendix B.

\begin{comment}
\subsection{Multi-Task Learning: Single-Task Experts vs Multi-Head}
\label{sec:multitask_analysis}

To directly test whether shared backbone improves generalization, we trained 8 independent expert networks (one per behavior) using CMA-ES, matching the multi-head training budget and protocol. Single-task experts achieved high training fitness ($0.95 \pm 0.04$ mean), yet failed catastrophically on generalization (Table~\ref{tab:singletask_vs_multihead}): GT1 score drops to $0.5597 \pm 0.3799$ (35\% degradation), with hard behaviors reaching only 8--30\% of their training performance. Multi-head architecture achieves consistent GT1 across all behaviors ($0.73$--$0.98$), demonstrating superior generalization. This pattern -- perfect single-task training but poor generalization -- aligns with multi-task learning theory (Caruana et al.~\cite{caruana1997multitask}), where shared backbone representations enable hard tasks to benefit from learning easier
tasks, thereby preventing overfitting.

\begin{table}[h]
\centering
\caption{Single-task expert networks vs multi-head CMA-ES: per-behavior GT1 generalization. Single-task networks achieve high training fitness but severe overfitting, especially on hard behaviors. Multi-head achieves balanced generalization across all 8 behaviors.}
\label{tab:singletask_vs_multihead}
\small
\begin{tabular}{lcccc}
\toprule
Behavior & Single-Task Train & Single-Task GT1 & Multi-Head GT1 & Gain \\
\midrule
\multicolumn{5}{l}{\textit{Easy-tier}} \\
stop & $1.000$ & $0.944$ & $0.949$ & $+0.005$ \\
motion\_reduction & $1.000$ & $0.936$ & $0.956$ & $+0.020$ \\
response\_magnitude & $1.000$ & $0.723$ & $0.980$ & $+0.257$ \\
elimination & $1.000$ & $0.731$ & $0.942$ & $+0.211$ \\
\midrule
\multicolumn{5}{l}{\textit{Hard-tier}} \\
motion\_increase & $0.958$ & $0.081$ & $0.731$ & $+0.650$ \\
vigor & $0.882$ & $0.139$ & $0.737$ & $+0.598$ \\
directedness & $0.850$ & $0.298$ & $0.732$ & $+0.434$ \\
confinement & $0.979$ & $0.627$ & $0.823$ & $+0.196$ \\
\midrule
Mean & $0.954$ & $0.560$ & $0.844$ & $+0.284$ \\
\bottomrule
\end{tabular}
\end{table}

\end{comment}

% -----------------------------------------------------------------------------------------------------------
\section{Discussion}
\label{sec:Discussion}

\subsection{What we learned}
\label{sec:learned}

Language-guided biological control is learnable from offline observational 
data without custom reward functions. The key architectural insight is that 
multi-task learning with multi-head decoupling and shared backbone 
substantially outperforms isolated single-task specialists (Appendix A.2): This pattern aligns with \cite{caruana1997multitask} -- shared representations enable hard tasks to benefit from easier ones, preventing 
overfitting. The network demonstrates robust semantic understanding, 
successfully generalizing to novel language inputs.

\subsection{Limitations}
\label{sec:limitations}

\textbf{Data scale and scope.} The archive is small (143 batches: 99 training, 44 test), potentially missing rare intervention patterns and limiting the range of duration-outcome relationships the network can learn. Evaluation covers only 8 behaviors in Xenobots; transfer to other organoid systems or whole organisms remains untested. The network selects from existing interventions rather than generating novel ones, restricting it to the parameter space explored during data collection.

\textbf{Reward signal validity.} The VLM reward is unvalidated against human expert raters, risking systematic bias or artifacts in the scoring pipeline. This is particularly concerning for subjective phenotypes (directedness, vigor) where human and algorithm disagreement could significantly skew results.

\textbf{Population specificity and the GT1-GT2 gap.} The sharp drop from 
GT1 (71.7\%) to GT2 (20.8\%) reveals a fundamental limitation: the network 
learns a population-specific mapping from prompt to intervention to phenotype. GT1 succeeds because test prompts pair with the same population of Xenobots. GT2 tests distribution shift: test videos come from a different biological population with distinct baseline activity, growth stage, and responsiveness. The network captures language semantics but not the underlying biological principles governing how interventions affect different populations.



\subsection{Future work}
\label{sec:futurework}

Our immediate priorities are human validation of VLM scoring, and active data collection targeting high-uncertainty regions of intervention space. Closing the loop with real-time experiments enables feedback-driven refinement. Architectural extensions include generative models to synthesize novel interventions, and lifelong learning for incremental behavior discovery. Transfer experiments to other biological systems would reveal whether language-to-intervention learning generalizes across cell types and scales.

% -----------------------------------------------------------------------------
\section*{Broader Impact}
\label{sec:broaderImpact}

This work translates natural-language descriptions into electrical 
interventions for Xenobots. The method lowers the barrier for researchers 
without electrophysiology expertise. A key risk is that non-experts may not 
recognize when an intervention has unintended consequences -- for example, a request to "make the bot disappear" retrieved an archived experiment where prolonged stimulation possibly destroyed the Xenobot. Since the system retrieves only archived interventions, this risk is bounded. Broader ethical questions about Xenobot research apply here as well.

% ------------------------------------------------------------------------------------------
%\section*{References}

\bibliographystyle{plain}  % or "unsrt", "apalike", etc.
\bibliography{neurips_2026}  % filename without .bib extension

\newpage
\appendix

\section*{Appendix}

\subsection*{Code and Data Availability}

Code, data, and analysis are provided on GitHub: \url{https://github.com/drartificialife/Neurips2026-Xenobots}. Upon publication acceptance, the Xenobots experimental data will be released according to institutional data policies and collaborator agreements.

\section*{Appendix A: P2I Network Architecture Details}
\label{app:architecture}

\subsection{Shared Backbone}
Input: 384-dimensional SentenceBERT embedding

Convolutional layers: Three 1D convolutions reducing dimensionality
\begin{align}
\text{Conv1D}(384 \to 64, k=3) &\to \text{ReLU} \\
\text{Conv1D}(64 \to 32, k=3) &\to \text{ReLU} \\
\text{Conv1D}(32 \to 16, k=3) &\to \text{ReLU}
\end{align}

Global average pooling: Output is $\mathbf{f}_{\text{shared}} \in \mathbb{R}^{16}$

\subsection{Per-Behavior Heads}
Each of 8 behavior heads shares the backbone but has independent downstream layers:

\begin{equation}
\mu_b = \sigma\left(\text{FC}_{32 \to 1}\left(\text{ReLU}\left(\text{FC}_{16 \to 32}(\mathbf{f}_{\text{shared}})\right)\right)\right), \quad b \in \{0, 1, \ldots, 7\}
\end{equation}

where $\sigma$ is sigmoid (output range $[0, 1]$ for intervention duration scaling).

\subsection{Parameter Count}

\paragraph{Backbone (shared Conv1d layers):}
\begin{align*}
\text{Conv1d}(384 \to 64, k=3) &: 384 \times 64 \times 3 + 64 = 73,792 \\
\text{Conv1d}(64 \to 32, k=3) &: 64 \times 32 \times 3 + 32 = 6,176 \\
\text{Conv1d}(32 \to 16, k=3) &: 32 \times 16 \times 3 + 16 = 1,552 \\
\text{Backbone Total} &: 73,792 + 6,176 + 1,552 = \mathbf{81,520 \text{ parameters}}
\end{align*}

\paragraph{Output heads (per behavior):}
\begin{align*}
\text{Linear}(16 \to 32) &: 16 \times 32 + 32 = 544 \\
\text{Linear}(32 \to 1) &: 32 \times 1 + 1 = 33 \\
\text{Per-head Total} &: 544 + 33 = \mathbf{577 \text{ parameters}}
\end{align*}

\paragraph{Network totals:}
\begin{align*}
\text{Single-head network} &: 7,984 + 577 = 8,561 \text{ parameters} \\
\text{Multi-head network (8 behaviors)} &: 7,984 + 8 \times 577 = 12,600 \text{ parameters} \\
\text{Overhead of multi-head} &: \frac{12,600 - 8,561}{8,561} = \frac{4,039}{8,561} = 47.2\%
\end{align*}

\subsection{Architecture Visualization}

\begin{figure}[h]
\centering
\includegraphics[width=0.95\textwidth]{figures/figure3.png}
\caption{P2I network architecture and VLM reward scoring. (A) Single-head baseline: one output for all 8 behaviors causes gradient interference on opposing objectives. Multi-head design: 8 independent heads (577 parameters each) with shared backbone (7,984 parameters) enable specialization without interference. Parameter overhead: 47.2\% (12,600 total for multi-head vs. 8,561 for single-head). (B) VLM reward scoring pipeline: pre/post motion heatmaps and prompt submitted to Qwen3.5-397B; VLM outputs score $s \in [0,1]$ for behavior-prompt alignment. Example (batch-000070): velocity increases 2.6×, yielding high scores for motion-increase behaviors and low scores for opposing behaviors. (C) VLM score distribution by behavior, showing easy-tier concentration (0.63--0.81) and hard-tier sparsity (0.14--0.19), justifying stratified train-test split and difficulty-based analysis.}
\label{fig:architecture}
\end{figure}

\subsection{Training Results Across 30 Seeds}

The multi-head architecture substantially outperforms the single-head baseline in training. Across 30 independent random seeds, the multi-head design achieves a mean final fitness of $0.927 \pm 0.010$ compared to single-head's $0.666 \pm 0.022$---a 39\% absolute improvement (Wilcoxon p = $7.54 \times 10^{-7}$, indicating highly significant superiority).

This advantage arises from the architectural decoupling: each behavior's head receives independent learning signals without interference from opposing objectives. The single-head baseline, by contrast, must compromise between conflicting goals (e.g., ``stop'' demands short stimulus duration while ``accelerate'' demands long duration), leading to poor optimization on the shared output.

Per-behavior analysis reveals that hard-tier behaviors (directedness, vigor, motion\_increase) benefit most from the multi-head design. These sparse-reward behaviors achieve dramatic improvements: directedness gains +0.60 (0.216 $\to$ 0.815), vigor +0.58 (0.254 $\to$ 0.833), and motion\_increase +0.25 (0.674 $\to$ 0.926). Easy-tier behaviors, already well-represented in the archive, show smaller absolute gains but consistently stable performance. This pattern reflects the regularization effect of multi-task learning: shared backbone features learned from easy behaviors help bootstrap learning on hard behaviors, preventing individual heads from overfitting to noise.

Figure~\ref{fig:training} visualizes these training dynamics across all 100 CMA-ES generations.

\begin{figure}[h]
\centering
\includegraphics[width=1\textwidth]{figures/training_combined.png}
\caption{\footnotesize Training results: (A) Overall training fitness curves comparing single-head (red dashed) vs. multi-head (green solid) architecture across 100 CMA-ES generations. Multi-head achieves 0.927 ± 0.010 final fitness vs. single-head 0.666 ± 0.022 (Wilcoxon p = 7.54e-07). (B) Per-behavior training fitness evolution for multi-head architecture, showing independent optimization trajectories for all 8 behaviors without gradient interference. (C) Boxplot of final per-behavior performance, highlighting dramatic improvements on hard behaviors (directedness +0.60, vigor +0.58, motion\_increase +0.25) while maintaining performance on easy behaviors.}
\label{fig:training}
\end{figure}

\section*{Appendix B: CMA-ES Hyperparameters and Configuration}
\label{app:cmaes}

\textbf{CMA-ES Variant:} We use the \textit{separable CMA-ES} (sep-CMA-ES) variant, which maintains a diagonal covariance matrix instead of full-rank. This reduces memory complexity from $O(n^2)$ to $O(n)$ and computational complexity from $O(n^3)$ to $O(n)$, making large-scale optimization tractable. For $n = 12,600$ parameters, sep-CMA-ES requires $\sim$50 MB memory (vs. $\sim$27 GB for full-rank CMA-ES), and scales linearly with parameter count.

Population size: $\lambda = 32$ (recommended: $\lambda = 4 + \lfloor 3 \ln(n) \rfloor = 4 + \lfloor 3 \times 9.44 \rfloor = 32$ with $n = 12,600$ total parameters; matches our actual choice)

Initial mutation strength: $\sigma_0 = 0.3$

Number of generations: 100 (training stops when fitness plateau is detected)

Combined fitness function:
\begin{equation}
\mathcal{F}_{\text{combined}} = \frac{1}{8} \sum_{i=1}^{8} \mathcal{F}_i(\mathbf{w})
\end{equation}

Each behavior's fitness is the mean VLM score across all training prompts assigned to that behavior.

\section*{Appendix C: Statistical Methods and Baseline Computation}
\label{app:statistical_methods}

\subsection{Wilcoxon Signed-Rank Test}

We use Wilcoxon signed-rank test (one-sided alternative: multi-head $>$ single-head) to compare final fitness across 30 seeds. This non-parametric test is appropriate for non-normally distributed fitness distributions.

\begin{align*}
H_0 &: \text{Multi-head fitness} = \text{Single-head fitness} \\
H_1 &: \text{Multi-head fitness} > \text{Single-head fitness}
\end{align*}

Significance level: $\alpha = 0.05$ (with Bonferroni correction for multiple behaviors: $\alpha/8 = 0.00625$)

\subsection{Effect Size}

Cohen's $d$ is computed as:
\begin{equation}
d = \frac{\mu_{\text{multi}} - \mu_{\text{single}}}{\sqrt{\frac{(n-1)s_{\text{multi}}^2 + (n-1)s_{\text{single}}^2}{2(n-1)}}}
\end{equation}

where $n=30$ seeds, $\mu$ is mean, $s$ is standard deviation.

\subsection{Baseline Computation}

Dataset baselines represent the expected VLM score from random batch selection without model guidance. For each behavior $i$:

\begin{equation}
\text{Baseline}_i = \frac{1}{N_{\text{prompts}} \times N_{\text{batches}}} \sum_{p \in P_i} \sum_{b \in B} S(b, p, i)
\end{equation}

where $S(b, p, i)$ is the VLM score for batch $b$, prompt $p$, and behavior $i$; $P_i$ is the set of training prompts for behavior $i$; $B$ is the set of training batches.

This gives per-behavior baseline scores representing the inherent difficulty: easy behaviors (high scores) have many good batches; hard behaviors (low scores) have few good batches.

\section*{Appendix D: Generalization Test Details}
\label{app:generalization_details}

\subsection{GT1: Generalization to Novel Prompts}
Test prompts: 5 held-out prompts per behavior (40 total, never seen during training)

Test batches: 99 training batches (same as training data)

Evaluation: For each of 40 test prompts, compute VLM score using the trained P2I network. Average score per behavior.

\subsection{GT2: Generalization to Novel Batches}
Test prompts: Same 5 held-out prompts per behavior

Test batches: 44 held-out test batches (never seen during training)

Evaluation: For each of 40 test prompts × 44 test batches combination, compute VLM score. Average score per behavior.

\section*{Appendix E: Prompt Lists}
\label{app:prompts}

\subsection{Training Prompts (40 total)}
\begin{verbatim}
stop: "stop moving", "remain stationary", "become immobile",
       "freeze in place", "cease all motion"

elimination: "make the bot disappeared", "bot ceases all activity",
             "complete motion shutdown", "bot becomes unresponsive",
             "permanent immobilization"

confinement: "stay confined to center", "localized activity",
             "restricted mobility", "confined movement",
             "center-seeking behavior"

motion_reduction: "move slower", "reduce movement", "minimize motion",
                  "decrease speed", "low velocity movement"

response_magnitude: "dramatic behavior change", "pronounced effect of intervention",
                    "stark difference from baseline", "strong behavior shift",
                    "large response to intervention"

directedness: "move in directed paths", "show purposeful movement",
              "goal-seeking behavior", "clear trajectory intent",
              "directed exploration"

vigor: "show vigor", "vigorous movement", "energetic behavior",
       "powerful movement", "high energy movement"

motion_increase: "move faster", "increase activity level",
                 "become more active", "high velocity movement",
                 "increase speed"
\end{verbatim}

\subsection{Test Prompts (40 total, novel variants)}
\begin{verbatim}
stop: "remain stationary", "become immobile", "freeze in place",
      "halt movement", "stop all motion"

elimination: "total cessation of movement", "bot viability lost",
             "lethal intervention effect", "irreversible behavioral change",
             "bot elimination"

confinement: "center-restricted movement", "confined activity",
             "limited mobility range", "centralized behavior",
             "avoid periphery"

motion_reduction: "sluggish motion", "lethargic behavior",
                  "minimize activity", "slower pace", "reduced speed"

response_magnitude: "substantial behavior shift", "marked behavioral change",
                    "pronounced difference from baseline", "significant effect",
                    "strong intervention response"

directedness: "intentional movement", "straight-line trajectory",
              "focused motion", "deliberate paths", "oriented behavior"

vigor: "forceful motion", "lively behavior", "robust movement",
       "intense activity", "spirited motion"

motion_increase: "accelerate motion", "boost activity", "rapid movement",
                 "increased velocity", "amp up speed"
\end{verbatim}

\section*{Appendix A.1: Multi-Task Learning Analysis}
\label{app:multitask}

\subsection{Single-Task Experts vs Multi-Head Architecture}

To directly test whether shared backbone improves generalization, we trained 8 independent expert networks (one per behavior) using CMA-ES, matching the multi-head training budget and protocol. Single-task experts achieved high training fitness ($0.95 \pm 0.04$ mean), yet failed catastrophically on generalization: GT1 score drops to $0.5597 \pm 0.3799$ (35\% degradation), with hard behaviors reaching only 8--30\% of their training performance. Multi-head architecture achieves consistent GT1 across all behaviors ($0.73$--$0.98$), demonstrating superior generalization.

This pattern -- perfect single-task training but poor generalization -- aligns with multi-task learning theory (Caruana et al.~\cite{caruana1997multitask}), where shared backbone representations enable hard tasks to benefit from learning easier tasks, thereby preventing overfitting.

\begin{table}[h]
\centering
\caption{Single-task expert networks vs multi-head CMA-ES: per-behavior GT1 generalization. Single-task networks achieve high training fitness but severe overfitting, especially on hard behaviors. Multi-head achieves balanced generalization across all 8 behaviors.}
\label{tab:singletask_vs_multihead}
\small
\begin{tabular}{lcccc}
\toprule
Behavior & Single-Task Train & Single-Task GT1 & Multi-Head GT1 & Gain \\
\midrule
\multicolumn{5}{l}{\textit{Easy-tier}} \\
stop & $1.000$ & $0.944$ & $0.949$ & $+0.005$ \\
motion\_reduction & $1.000$ & $0.936$ & $0.956$ & $+0.020$ \\
response\_magnitude & $1.000$ & $0.723$ & $0.980$ & $+0.257$ \\
elimination & $1.000$ & $0.731$ & $0.942$ & $+0.211$ \\
\midrule
\multicolumn{5}{l}{\textit{Hard-tier}} \\
motion\_increase & $0.958$ & $0.081$ & $0.731$ & $+0.650$ \\
vigor & $0.882$ & $0.139$ & $0.737$ & $+0.598$ \\
directedness & $0.850$ & $0.298$ & $0.732$ & $+0.434$ \\
confinement & $0.979$ & $0.627$ & $0.823$ & $+0.196$ \\
\midrule
Mean & $0.954$ & $0.560$ & $0.844$ & $+0.284$ \\
\bottomrule
\end{tabular}
\end{table}

\subsection{Mechanisms of Multi-Task Learning Success}

Following Caruana et al., we identify three mechanisms enabling multi-head success:

\textit{Representation bias toward shared minima.} The shared backbone simultaneously optimizes 8 behaviors, biasing search toward representations useful to all tasks rather than optimal for any single task. Single-task networks, trained in isolation, optimize separately and often converge to local minima that overfit to task-specific noise. Hard behaviors (motion\_increase, directedness, vigor) suffer most, achieving only 8--30\% generalization despite perfect training. Shared learning forces discovery of robust features that work across behaviors.

\textit{Eavesdropping: hard tasks learn from easy tasks.} Easy behaviors (stop, motion\_reduction) provide strong, stable gradient signals. Hard behaviors, with sparser reward landscapes, can ``listen to'' feature learning from easy tasks through the shared backbone. For example, the clear motion signal from stop learning helps bootstrap representations that enable directedness learning. Single-task networks lose this auxiliary signal, leaving hard behaviors to learn from weak gradients alone.

\textit{Automatic regularization via multi-objective balance.} Training 8 objectives simultaneously prevents overfitting to spurious features in any single task. Multi-head's slightly lower training fitness ($0.876$ vs $0.954$ single-task) reflects necessary compromises; yet this produces dramatically higher generalization ($0.844$ GT1 vs $0.560$ single-task), indicating better regularization. The network learns to prioritize features generalizable across behaviors.

The data strongly suggest that for this class of problems---learning from sparse, noisy signals across multiple related objectives---multi-task learning with shared backbone is not optional but essential.

\section*{Appendix F: Per-Behavior Generalization Statistics (GT1 and GT2)}
\label{app:per_behavior_stats}

This appendix provides detailed per-behavior breakdowns of generalization test performance.

\subsection{Generalization to Novel Prompts (GT1)}

\begin{table}[h]
\centering
\caption{GT1 per-behavior statistics: novel test prompts paired with training-set videos. Computed across 30 seeds; $p$-values from Wilcoxon signed-rank tests against the random-retrieval baseline.}
\label{tab:gt1_detailed}
\small
\begin{tabular}{lcccccc}
\toprule
Behavior & Baseline & GT1 mean & Improvement & $p$-value & Cohen's $d$ & $R^2$ \\
\midrule
\texttt{stop}                & 0.7496 & 0.9517 & $+27.0\%$  & $6.80 \times 10^{-8}$ & 22.52 & 0.998 \\
\texttt{motion\_reduction}   & 0.7633 & 0.9030 & $+18.3\%$  & $2.09 \times 10^{-7}$ & 2.58 & 0.870 \\
\texttt{motion\_increase}    & 0.1397 & 0.7340 & $+425.3\%$ & $7.05 \times 10^{-7}$ & 9.23 & 0.988 \\
\texttt{vigor}               & 0.1559 & 0.7410 & $+375.3\%$ & $4.17 \times 10^{-7}$ & 20.00 & 0.998 \\
\texttt{response\_magnitude} & 0.7623 & 0.9800 & $+28.6\%$  & $1.01 \times 10^{-7}$ & 42.15 & 0.999 \\
\texttt{directedness}        & 0.1502 & 0.7323 & $+387.6\%$ & $2.61 \times 10^{-7}$ & 11.82 & 0.993 \\
\texttt{confinement}         & 0.6309 & 0.9443 & $+49.7\%$  & $4.73 \times 10^{-7}$ & 2.32 & 0.843 \\
\texttt{elimination}         & 0.6995 & 0.9383 & $+34.1\%$  & $3.32 \times 10^{-7}$ & 6.89 & 0.979 \\
\bottomrule
\end{tabular}
\end{table}

All eight behaviors clear the $p < 10^{-7}$ threshold with Cohen's $d$ ranging from 2.32 to 42.15, indicating large effect sizes. Variance explained ($R^2$) ranges from 0.84 to 0.999, indicating that the learned mapping accounts for most score variation across novel prompts.

\subsection{Generalization to Novel Batches (GT2)}

\begin{table}[h]
\centering
\caption{GT2 per-behavior statistics: novel test prompts paired with held-out test videos. Computed across 30 seeds; $p$-values from Wilcoxon signed-rank tests against the random-retrieval baseline. Negative improvements indicate scores below baseline.}
\label{tab:gt2_detailed}
\small
\begin{tabular}{lcccccc}
\toprule
Behavior & Baseline & GT2 mean & Improvement & $p$-value & Cohen's $d$ & $R^2$ \\
\midrule
\texttt{stop}                & 0.7496 & 0.8717 & $+16.3\%$  & $8.66 \times 10^{-3}$ & 0.37 & 0.123 \\
\texttt{motion\_reduction}   & 0.7633 & 0.5580 & $-26.9\%$  & $6.34 \times 10^{-7}$ & $-1.08$ & 0.537 \\
\texttt{motion\_increase}    & 0.1397 & 0.3993 & $+185.8\%$ & $9.74 \times 10^{-7}$ & 1.44 & 0.675 \\
\texttt{vigor}               & 0.1559 & 0.6337 & $+306.5\%$ & $4.17 \times 10^{-7}$ & 4.97 & 0.961 \\
\texttt{response\_magnitude} & 0.7623 & 0.9783 & $+28.3\%$  & $1.30 \times 10^{-7}$ & 2.48 & 0.860 \\
\texttt{directedness}        & 0.1502 & 0.3420 & $+127.7\%$ & $6.02 \times 10^{-7}$ & 2.77 & 0.885 \\
\texttt{confinement}         & 0.6309 & 0.2960 & $-53.1\%$  & $9.87 \times 10^{-5}$ & $-0.87$ & 0.433 \\
\texttt{elimination}         & 0.6995 & 0.9357 & $+33.8\%$  & $7.49 \times 10^{-7}$ & 2.25 & 0.835 \\
\bottomrule
\end{tabular}
\end{table}

Six of eight behaviors show significant improvement over baseline ($p < 10^{-5}$). Two behaviors --- \texttt{confinement} and \texttt{motion\_reduction} --- show statistically significant negative transfer ($p < 10^{-4}$), likely due to distribution shift in test cohort duration statistics. The antagonistic behaviors (\texttt{vigor}, \texttt{motion\_increase}, \texttt{directedness}) maintain substantial gains ($+127\%$ to $+307\%$) even under the harder GT2 protocol.

\subsection{Correlation Analysis: GT1 vs GT2}

Pearson correlation between GT1 and GT2 performance per behavior:

\begin{table}[h]
\centering
\caption{Pearson correlation: GT1 vs GT2 fitness across 30 seeds.}
\begin{tabular}{lcc}
\toprule
Behavior & Pearson $r$ & $p$-value \\
\midrule
stop & 0.067 & 0.724 \\
elimination & 0.029 & 0.879 \\
confinement & -0.487 & 0.006 ** \\
motion\_reduction & 0.094 & 0.620 \\
response\_magnitude & 0.370 & 0.044 * \\
directedness & 0.660 & $<0.001$ *** \\
vigor & -0.391 & 0.033 * \\
motion\_increase & -0.149 & 0.433 \\
\bottomrule
\multicolumn{3}{l}{\small *** $p < 0.001$, ** $p < 0.01$, * $p < 0.05$}
\end{tabular}
\end{table}

Weak correlations suggest that GT1 and GT2 measure different generalization modes: semantic generalization (novel language) vs. biological transfer (novel batches). Some behaviors (directedness, response\_magnitude) show positive correlation (good GT1 $\to$ good GT2), while others show negative correlation (good GT1 but poor GT2 on novel batches), indicating archive-limited behaviors.

\section*{Appendix G: REINFORCE vs CMA-ES Ablation Study}
\label{app:reinforce_ablation}

We compared CMA-ES with a policy gradient baseline (REINFORCE with baseline) to understand whether the superior CMA-ES performance reflects the choice of optimization algorithm or simply the need for more data/computation. This section documents the ablation.

\subsection{Why REINFORCE for Contextual Bandits?}

Our problem is a contextual bandit: the environment provides a state (prompt embedding), the agent selects an action (intervention duration), and receives a scalar reward (VLM score) with no state transition. REINFORCE~\cite{williams1992simple} is theoretically justified for this setting because:

\begin{itemize}
\item \textbf{Terminal rewards only:} No trajectory to optimize; single sample per episode suffices.
\item \textbf{Fixed state representation:} Prompts are embedded to fixed 384-dim vectors; the policy learns a deterministic mapping.
\item \textbf{Single action and reward:} No need for temporal credit assignment; advantage estimation reduces to reward minus baseline.
\item \textbf{Sparse, noisy rewards:} Running-mean baseline \cite{williams1992simple} provides automatic variance reduction without learning a separate value network.
\end{itemize}

\subsection{REINFORCE Policy Gradient Formulation}

The policy $\pi_\theta$ maps prompt embedding $\mathbf{e} \in \mathbb{R}^{384}$ and behavior $b$ to a scalar action $a \in [0, 1]$ representing predicted intervention duration:

\begin{equation}
a = \pi_\theta(\mathbf{e}, b) = \sigma(\text{MLP}(\mathbf{e}; \Theta_b))
\end{equation}

where $\sigma$ is sigmoid and $\Theta_b$ are behavior-specific parameters.

The policy gradient objective is derived from the REINFORCE algorithm \cite{williams1992simple}:
\begin{equation}
\nabla_\theta J(\theta) = \mathbb{E}_{b,\mathbf{e}}[(G_t - V_b) \nabla_\theta \log \pi_\theta(a \mid \mathbf{e}, b)]
\end{equation}

where $G_t = r$ is the terminal VLM reward and $V_b = \text{RunningMean}(R_{\text{history}}(b))$ is a per-behavior baseline for variance reduction. The running mean is updated incrementally:

\begin{equation}
V_b^{(k+1)} = \frac{k \cdot V_b^{(k)} + r_k}{k+1}
\end{equation}

The training loss includes entropy regularization to encourage exploration:
\begin{equation}
\mathcal{L}(\theta) = -\mathbb{E}[(r - V_b) \log \pi_\theta(a \mid \mathbf{e}, b)] + \beta H(\pi_\theta)
\end{equation}

where $\beta = 0.01$ and $H(\pi_\theta) = -\sum_a \pi_\theta(a \mid \mathbf{e}, b) \log \pi_\theta(a \mid \mathbf{e}, b)$ is the policy entropy.

\subsection{REINFORCE Algorithm}

\begin{algorithm}[h]
\caption{REINFORCE for P2I Network Training}
\begin{algorithmic}[1]
\State \textbf{Input:} Network $\pi_\theta$, learning rate $\alpha = 10^{-3}$, entropy weight $\beta = 0.01$
\State \textbf{Initialize:} Random weights $\theta_0$, baseline history $H_b = \{\}$ for each behavior $b \in \{1, \ldots, 8\}$
\For{episode $k = 1, \ldots, 76800$}
  \State Sample behavior: $b \sim \text{Uniform}(\{1, \ldots, 8\})$
  \State Sample prompt: $p \sim \text{Dataset}(b)$
  \State Encode: $\mathbf{e} = \text{SentenceBERT}(p) \in \mathbb{R}^{384}$
  \State Forward pass: $a = \pi_\theta(\mathbf{e}, b)$  (\textit{scalar action in [0,1]})
  \State Map to duration: $d_{\text{pred}} = d_{\min} + a \cdot (d_{\max} - d_{\min})$
  \State Find nearest batch: $b_{\text{near}} = \arg\min_j |d_j - d_{\text{pred}}|$
  \State Retrieve reward: $r = \text{VLM}(b, p, b_{\text{near}})$ from precomputed scores
  \State Update baseline: $H_b \gets H_b \cup \{r\}$
  \State Compute advantage: $A = r - \text{Mean}(H_b)$
  \State Compute loss: $\mathcal{L} = -A \log \pi_\theta(a \mid \mathbf{e}, b) + \beta H(\pi_\theta(\cdot \mid \mathbf{e}, b))$
  \State Gradient step: $\theta \gets \theta + \alpha \nabla_\theta \mathcal{L}$
\EndFor
\State \textbf{Output:} Trained policy $\pi_\theta$
\end{algorithmic}
\end{algorithm}

\subsection{Training Configuration}

REINFORCE training was configured to match CMA-ES' computational budget:

\begin{table}[h]
\centering
\caption{REINFORCE vs CMA-ES computational alignment.}
\begin{tabular}{lcc}
\toprule
\textbf{Metric} & \textbf{REINFORCE} & \textbf{CMA-ES} \\
\midrule
Episodes / Generations & 76,800 & 100 \\
Network evaluations per step & 1 & 32 \\
Behavior-prompt pairs per evaluation & 1 & $8 \times 3 = 24$ \\
\midrule
\textbf{Total behavior-prompt pairs} & \textbf{76,800} & \textbf{76,800} \\
\midrule
Learning rate & 0.001 & — \\
Entropy weight & 0.01 & — \\
Baseline & Running mean/behavior & — \\
\bottomrule
\end{tabular}
\end{table}

\subsection{Computational Budget Alignment}

Both methods evaluate the same number of (behavior, prompt) pairs. CMA-ES performs:
\begin{equation}
100 \text{ generations} \times 32 \text{ population} \times 8 \text{ behaviors} \times 3 \text{ prompts/behavior} = 76,800 \text{ VLM evaluations}
\end{equation}

REINFORCE samples behavior-prompt pairs uniformly during training:
\begin{equation}
76,800 \text{ episodes} \times 1 \text{ (behavior, prompt) pair/episode} = 76,800 \text{ VLM evaluations}
\end{equation}

Thus both methods have equal computational cost in terms of VLM calls, ensuring a fair comparison.

\subsection{Per-Behavior Results}

Table~\ref{tab:reinforce_per_behavior} shows per-behavior training fitness across CMA-ES single-head, CMA-ES multi-head, and REINFORCE:

\begin{table}[h]
\centering
\caption{Per-behavior training fitness: CMA-ES vs REINFORCE. CMA-ES: final generation, mean across 30 seeds. REINFORCE: final epoch, single run. Baseline VLM scores are computed on untrained random network.}
\label{tab:reinforce_per_behavior}
\begin{tabular}{lcccccc}
\toprule
Behavior & Baseline & Single-Head & Multi-Head & REINFORCE & CMA-ES/RF Ratio \\
\midrule
\texttt{stop}                & 0.748 & 0.937 & 0.955 & 0.634 & 1.51 \\
\texttt{motion\_reduction}   & 0.763 & 0.872 & 0.967 & 0.809 & 1.19 \\
\texttt{motion\_increase}    & 0.140 & 0.674 & 0.926 & 0.338 & 2.74 \\
\texttt{vigor}               & 0.156 & 0.254 & 0.833 & 0.285 & 2.92 \\
\texttt{response\_magnitude} & 0.768 & 0.958 & 0.997 & 0.820 & 1.22 \\
\texttt{directedness}        & 0.148 & 0.216 & 0.815 & 0.190 & 4.29 \\
\texttt{confinement}         & 0.648 & 0.756 & 0.967 & 0.640 & 1.51 \\
\texttt{elimination}         & 0.702 & 0.666 & 0.955 & 0.545 & 1.75 \\
\bottomrule
\end{tabular}
\end{table}

\subsection{Analysis: CMA-ES vs REINFORCE Performance and Mechanisms}

\textbf{Empirical Performance Disparities.} Hard-tier behaviors reveal the largest gaps between methods. REINFORCE fails catastrophically on sparse-reward problems:
\begin{itemize}
  \item \texttt{Directedness}: REINFORCE 0.190 vs CMA-ES 0.815 (4.3$\times$ gap)
  \item \texttt{Vigor}: REINFORCE 0.285 vs CMA-ES 0.833 (2.9$\times$ gap)
  \item \texttt{Motion\_increase}: REINFORCE 0.338 vs CMA-ES 0.926 (2.7$\times$ gap)
\end{itemize}
The policy gradient method barely improves over baseline (only +28\% on directedness), while CMA-ES achieves +452\%. Easy behaviors show mixed results; REINFORCE occasionally matches CMA-ES on response\_magnitude (0.820 vs 0.997, 1.22$\times$ gap) and motion\_reduction (0.809 vs 0.967, 1.19$\times$ gap), but consistently underperforms. Averaging across all behaviors: REINFORCE 0.533 vs CMA-ES multi-head 0.927 (73\% improvement).

\textbf{Why REINFORCE Fails: Three Mechanisms.} The policy gradient method struggles for fundamental reasons:

\begin{enumerate}
  \item \textbf{High variance on sparse rewards.} REINFORCE uses single-trajectory gradient estimates with variance scaling as $O(T^2)$ in episode length. With only 9,600 samples per behavior (76,800 episodes $\div$ 8 behaviors), the gradient signal is drowned by noise. Even with baseline subtraction, variance remains $O(T)$, insufficient for convergence on behaviors where $P(\text{good intervention}) \approx 0.3$.

  \item \textbf{Poor credit assignment on multi-head network.} The network has 8 independent output heads, one per behavior. During training on behavior $b_j$, only head $j$ receives learning signals. The shared backbone simultaneously receives conflicting gradients from all 8 behaviors, each pulling the shared weights toward different optima. CMA-ES avoids this by optimizing the entire weight vector as a unit, naturally balancing multi-objective pressure.

  \item \textbf{Unreliable gradient estimation on discrete behavior selection.} Although the network outputs a continuous duration value, it implicitly selects a behavior during sampling. The gradient estimate $\nabla_\theta \log \pi_\theta(a|e,b) \cdot (r - V_b)$ is biased when the discrete behavior choice has not been explored sufficiently; baseline subtraction cannot correct this discrete-action bias.
\end{enumerate}

CMA-ES avoids these issues through population-based exploration and black-box optimization, making it more robust on sparse-reward, high-dimensional, multi-objective problems.

\section*{Appendix H: Video Processing and Prompt Routing}
\label{app:video_and_routing}

\subsection{Video Processing: Motion Denoising Pipeline}

Raw xenobot videos contain optical noise (bubbles, electrical reflections) that confounds motion analysis. We denoise using a four-step pipeline:

\begin{enumerate}
\item \textbf{Frame differencing:} Consecutive frames are differenced to detect motion regions.
\item \textbf{HSV color thresholding:} White pixels (low saturation $< 20$, corresponding to bubbles and reflections) are rejected; colored pixels (high saturation, corresponding to pigmented xenobots) are retained.
\item \textbf{Morphological operations:} Erosion and dilation clean isolated noise artifacts.
\item \textbf{Optical flow on masked regions:} Dense optical flow (Farneback method) is computed only on the bot mask, yielding motion magnitude heatmaps.
\end{enumerate}

The motion heatmap is computed as:
\begin{equation}
\Delta H = \sum_t \text{Magnitude}_{\text{masked}}^{\text{post}}(t) - \sum_t \text{Magnitude}_{\text{masked}}^{\text{pre}}(t)
\end{equation}
where magnitude is accumulated across all consecutive frame pairs in the pre- and post-intervention periods.

\begin{figure}[h]
\centering
\includegraphics[width=0.8\textwidth]{figures/figure1.png}
\caption{\footnotesize Motion denoising pipeline for xenobot behavioral analysis.
(a) Raw video frame showing multi-electrode array chamber with xenobot. Bubbles
(translucent white spheres) and electrical reflections create optical noise.
(b) Optical flow magnitude heatmap computed from raw consecutive frames (frame $t$ to
$t+5$). Red indicates motion regions; extensive noise from bubbles and reflections
throughout the chamber confounds true xenobot motion signal.
(c) Frame differencing mask: bright regions indicate pixels with intensity change between
frames. Detects all motion (xenobot and bubbles).
(d) HSV saturation mask: white pixels (low saturation $< 20$) corresponding to bubbles
and reflections are rejected; colored pixels (high saturation) corresponding to pigmented
xenobots are retained.
(e) Combined bot mask: element-wise intersection of motion and color masks. Xenobot
regions preserved; bubble noise eliminated through color filtering.
(f) Cleaned optical flow magnitude heatmap: only motion in bot regions shown. This
represents true xenobot behavior without artifact noise. Magnitude values accumulated
across all consecutive frame pairs to compute $\Delta H$, the net behavioral change
between pre- and post-intervention periods.}
\label{fig:denoising}
\end{figure}

\subsection{Prompt Embeddings and Routing}

Prompts are embedded using SentenceBERT (SBERT), a sentence transformer pretrained on natural language inference tasks. SBERT produces 384-dimensional embeddings that capture semantic similarity.

\subsubsection{Training Routing (Oracle)}
During training, prompts are assigned to their true behavior class (oracle routing). The CMA-ES optimizer then maximizes fitness for each behavior's assigned head.

\subsubsection{Test Routing (Cosine Similarity)}
At test time, the true behavior label is not available. Instead, we use cosine similarity routing: each test prompt is routed to the training prompt (from the 40 training prompts) with highest embedding similarity. The model then uses the corresponding behavior head.

This approach allows generalization to novel prompts without requiring the model to learn behavior classification; it leverages the assumption that semantically similar prompts should map to the same behavior head.

\newpage
\input{checklist.tex}



\end{document}